% \subsubsection{An Example}
% \label{sec:worked_example_qft_qpe}

% This subsection gives a concrete example of (i) sampling an instance-level configuration,
% (ii) instantiating two generators with explicit local parameters, (iii) merging the resulting
% subcircuits into a final circuit, and (iv) extracting static, graph, and dynamic features.

% \paragraph{Step 1: Draw an instance-level configuration.}
% For a single benchmark instance, we sample
% \begin{equation}
% (n_q, d, k, r, \xi),
% \end{equation}
% as defined in Sec.~3.1.3, where $n_q$ is qubits, $d$ is a target depth budget,
% $k$ is the evaluation-qubit budget, $r$ is a repetition count for repeatable sections,
% and $\xi$ collects generator-specific configuration variables (e.g., QFT swap flags). %
% \;\;\;%
% (Example draw) we set
% \begin{equation}
% n_q = 7,\quad d = 80,\quad k = 4,\quad r = 1.
% \end{equation}
% We index the qubits as $Q=\{q_0,\dots,q_6\}$.

% \paragraph{Step 2: Instantiate generator \#1 (QFT) with explicit local parameters.}
% Let the first generator be a QFT acting on a 3-qubit system register
% \begin{equation}
% Q_S = \{q_0,q_1,q_2\}.
% \end{equation}
% We draw generator-local settings
% \begin{equation}
% \begin{aligned}
% \xi_{\mathrm{QFT}} 
% &= (\texttt{include\_swaps}=\texttt{false},  \texttt{inverse}=\texttt{false}, \\
% &\quad \texttt{entangled}=\texttt{false}).
% \end{aligned}
% \end{equation}
% This yields a realized subcircuit
% \begin{equation}
% S_1 := C_{\mathrm{QFT}}(Q_S;\xi_{\mathrm{QFT}}),
% \end{equation}
% where the structural template (qubit layout, gate topology, depth contribution) is fixed
% by the generator and $Q_S$ (Sec.~3.1.3). 

% \paragraph{Step 3: Instantiate generator \#2 (QPE) with explicit local parameters.}
% Let the second generator be QPE, using a 4-qubit phase register and the same system register:
% \begin{equation}
% Q_P = \{q_3,q_4,q_5,q_6\},\qquad Q_S = \{q_0,q_1,q_2\}.
% \end{equation}
% We draw QPE-local settings
% \begin{equation}
% \xi_{\mathrm{QPE}} = (m=|Q_P|=4,\; \texttt{iqft\_include\_swaps}=\texttt{false},\; U = R_Z(\theta)\text{ on }q_0),
% \end{equation}
% with an explicit sampled gate parameter
% \begin{equation}
% \theta = 1.998\;\text{rad}.
% \end{equation}
% (Concrete realization.) We prepare an eigenstate of $U$ by applying $X$ on $q_0$ (so $q_0$ is in $|1\rangle$),
% apply Hadamards on the phase register, then apply the controlled powers
% \begin{equation}
% \text{c-}U^{2^j} = \text{c-}R_Z(2^j\theta)\quad \text{from control } q_{3+j}\ \text{to target } q_0,\ \ j\in\{0,1,2,3\}.
% \end{equation}
% For this example, the numerical angles are
% \begin{align}
% 2^0\theta &= 1.998\ \text{rad}, &
% 2^1\theta &= 3.996\ \text{rad}, \\
% 2^2\theta &= 7.992\ \text{rad}, &
% 2^3\theta &= 15.984\ \text{rad},
% \end{align}
% (where rotations are understood modulo $2\pi$). This yields
% \begin{equation}
% S_2 := C_{\mathrm{QPE}}(Q_P,Q_S;\xi_{\mathrm{QPE}}).
% \end{equation}
% This example mirrors Sec.~3.1.3: a structural template is instantiated first, and numerical
% gate values (here $\theta$) are then sampled as generator-specific parameters. 

% \paragraph{Step 4: Merge into the final circuit.}
% After generator realization, the final benchmark circuit is the ordered composition
% \begin{equation}
% C = S_1 \circ S_2,
% \end{equation}
% as in Sec.~3.1.5. 

% \paragraph{Step 5: Extract features for the same circuit.}
% Following Sec.~3.2, the circuit record stores three feature classes. 

% \textbf{Static metrics} are computed from the circuit description/metadata (Table~2),
% e.g., \path|num_qubits|, \path|width|, \path|depth|, \path|circuit_size|,
% \path|two_qubit_gate_count|, \path|two_qubit_gate_percentage|, and \path|gate_counts|. 

% \textbf{Graph metrics} are computed from the qubit interaction graph where vertices are qubits
% and edges represent multi-qubit interactions (Sec.~3.2.2).
% In this example, the controlled-$R_Z$ operations in QPE induce edges between each phase qubit
% $q_{3+j}$ and the target $q_0$, increasing degree-based descriptors such as \path|max_degree|
% and affecting path-length descriptors such as \path|diameter| and \path|average_shortest_path_length|
% (Table~3). 

% \textbf{Dynamic metrics} depend on the output state and require strong simulation; \sys includes
% \path|statevector_saved_entropy| and \path|statevector_saved_sparsity| with exponential $O(2^N)$
% time and space cost (Table~4). 

 % Preamble (once)
 







\subsection{Running Example}
\label{sec:worked_example_qft_qpe}

Figure~\ref{fig:sec41_running_example} illustrates one end-to-end circuit instance generated by \sys.
The workflow is intentionally analogous to a database benchmark that (i)  selects a template sequence, (ii) samples a scale/configuration,
(iii) instantiates a small sequence of templates with parameters, and  (iv) materializes both an artifact and its metadata.

\textbf{Step 1 (Template selection).} Starting from START, \sys samples a short
template sequence using the history-dependent distribution. Here the sequence is
\emph{QFT} $\rightarrow$ \emph{QPE}.

\textbf{Step 2--3 (Parameters \& subcircuits generation).}  Given a user seed, \sys samples an instance configuration
$(n_q, d, k, r, \xi)$. In this example, $n_q{=}7$ qubits, depth  $d{=}80$, evaluation  $k{=}4$,
and repetition count $r{=}1$. Each template is then instantiated by binding
template-local parameters (e.g., boolean flags and numeric values such as $\phi{=}1.998$). All choices are deterministic given the seed.

\textbf{Step 4 (Merge \& provenance).} The instantiated subcircuits are composed in order to produce the
final circuit artifact $C = S_1 \circ S_2$. Alongside $C$, \sys stores a provenance record containing the
template sequence and all sampled parameters, enabling exact replay and database-style inspection.

% \textbf{Step 5 (Metrics).} From the same circuit $C$, \sys extracts three metric classes:
% (i) \emph{static} metrics (e.g., qubit count, depth, gate counts) computed by scanning the circuit;
% (ii) \emph{graph} metrics computed from the induced interaction graph (two-qubit operations create edges);
% and (iii) \emph{dynamic} metrics (e.g., entropy/sparsity) computed via simulation, which is substantially more expensive.
